\documentclass[10pt]{article}
\usepackage[a4paper,left=2.54cm,top=2.54cm,right=2.54cm,bottom=2.54cm]{geometry}
\usepackage{fancyhdr}
\setlength{\headsep}{1.40cm} % Adjust the space after the header
\usepackage{afterpage}
\usepackage{setspace}
\usepackage{iccfdTemplate/bibspacing} 
\singlespacing

%%%% YOU CAN PUT YOUR OWN DEFINITIONS HERE
\newfont{\toto}{msbm10 at 12 pt}
\newfont{\ithd}{cmr9}

\newcommand{\si}[1]{\rm\scriptscriptstyle{#1}}
%%%% END OF YOUR DEFINITIONS 
 
\pagestyle{fancyplain}
\renewcommand{\headrulewidth}{0pt}

\usepackage{amsmath,amsthm,amsfonts,amssymb,mathtools}
\usepackage[pdftex]{graphicx} 
\usepackage[T1]{fontenc}

%%! Additional usepackage
\usepackage{physics}
\usepackage[hidelinks]{hyperref}
% \usepackage{hyperref} 
% \hypersetup{
	%     colorlinks=true,
	%     linkcolor=blue,
	%     filecolor=magenta,       
	%     urlcolor=cyan, 
	%     } 
\usepackage{subcaption}

\usepackage{csquotes}
\usepackage{enumitem}

%%%% CONFERENCE HEADER. REPLACE xxxx WITH 4-DIGIT PAPER NUMBER ASSIGNED BY CONFERENCE COMMITTEE.

% \rhead{\ithd{\bf ICCFD12-2024-xxxx\\  \   \\}}
\lhead{\ithd{\bf Twelfth International Conference on \\      
Computational Fluid Dynamics (ICCFD12), \\
Kobe, Japan, July 14-19, 2024
}}
% Define the left and center headers to be empty or contain fixed content
\rhead{}
\chead{}

\usepackage{titling}
\setlength{\droptitle}{0em}  
\pretitle{\vspace{-4em}\begin{center}\LARGE}
\posttitle{\end{center}\vspace{-1em}}
\preauthor{\begin{center}\large}
\postauthor{\end{center}\vspace{-3em}}


\title{
\bf Positivity-preserving algorithm for implicit finite
volume methods simulating compressible flows
}
\author{
Qian Wang$^{*}$, Hanyu Zhou$^{**}$\\
Corresponding author: qian.wang@csrc.ac.cn \\
$^{*}$ Beijing Computational Science Research Center, Beijing, 100193, China.\\
$^{**}$ Department of Engineering Mechanices, Tsinghua Universiy, Beijing, 100084, China.
}
\date{}

\begin{document}

%%%% TITLE
\maketitle
\afterpage{\fancyhead{}}
%\pagenumbering{gobble}

%%%% ABSTRACT AND KEYWORDS
%\vskip0.5cm
\centerline{
\begin{minipage}[t]{150mm}
{\bf Abstract:} 
A commonly encountered issue in the numerical simulation of challenging compressible flow problems is the failure to preserve positivity of density or pressure. 
%The majority of existing positivity-preserving algorithms are designed for numerical schemes based on explicit time stepping. 
A positivity-preserving algorithm for finite volume methods using dual-time stepping is developed in this paper, towards robust numerical simulations of compressible flows. In the positivity-preserving algorithm, admissible cell-averages are obtained by controlling physical and pseudo time step sizes which results in limited solution changes. To overcome the difficulty of unknown solution changes in time step limiting, we propose to use explicit time discretizations to obtain efficient estimations of future states. The allowable time step sizes are determined by limiting the relative solution changes to ensure positivity of updated solutions. Given positive cell-averages, admissible reconstruction polynomials can be obtained by applying a positivity-preserving scaling limiter. The positivity-preserving algorithm is applied to a high-order variational finite volume method with an explicit first stage singly diagonally implicit Runge-Kutta (ESDIRK) time integration. Numerical results for a series of benchmark test cases demonstrate the robustness and high resolution of the positivity-preserving implicit high-order finite volume method. 

\vskip0.2cm
{\it Keywords:} Positivity Preserving, Finite Volume, Implicit Time Stepping, Compressible Flows. \\
\end{minipage}
}
\vskip0.5cm

%%%% MAIN PART
% \section{Introduction}
% This is the main part of the paper. Its length should be no longer than {\bf 25 pages}, including key figures and references. 
% It can be divided in as many sections as you decide. The paper must be prepared using this template 
% and compiled using standard \LaTeX, generating a {\bf PDF file} that will be finally uploaded. If
% another kind of word processor is utilized, please adhere to the formatting provided in the PDF template.


% \section{Problem Statement}
% This document allows you to easily include references \cite{book,journalpaper}, equations, figures (see Figure \ref{fig:logo}) or anything else you
% desire into a clean and compact environment of \LaTeX.  For example if you'd like to impress a date you can write
% the unsteady heat equation as
% \begin{eqnarray}
% \frac{\partial \mathbf{V}}{\partial t} - \alpha \left( \frac{\partial^2 \mathbf{V}}{\partial x^2} +
%        \frac{\partial^2 \mathbf{V}}{\partial y^2} +
%        \frac{\partial^2 \mathbf{V}}{\partial z^2} \right)
% = 0
% \label{eq:heat}
% \end{eqnarray}
% where $x, y, z$ are the space dimensions and $\alpha$ is a parameter.  If you felt inclined you could define $\mathbf{V}$ as
% %
% $$\mathbf{V} = y^2 z - \text{cos}(0.1 x)$$
% %
% for a non-exact solution.  Computational fluid dynamics~\cite{paper} can be used to discretize the equations, apply boundary conditions and 
% simulation the unsteady nature of the flow.  An innovative method to simulate the heat equation could even be submitted to ICCFD12.

%   The scope of ICCFD12~\cite{web} is devoted to all innovative aspects of CFD, basic and applied.  Subjects of interest include but are not limited
% to: 
% \begin{itemize}

%     \item Innovative algorithm development for flow simulations: higher-order methods, iterative methods, parallel
%           algorithms, mesh adaptation, grid generation, meshless methods, immersed boundary methods and level-set
%           methods. 
          
%     \item Advances in modeling of flow physics in the area of: steady and unsteady flows, compressible and incompressible 
%           flows, flows in porous media, hypersonic and reacting flows, turbulence (transition, DNS/LES, etc.), multi-phase 
%           flows, boundary layer stability and vortex dynamics. 
          
%     \item Advanced multidisciplinary applications using the above mentioned technologies: aeroacoustics, flow
%           control, biomedical fluid mechanics, large scale applications, verification and validation methods, and
%           turbomachinery. 
          
% \end{itemize}

% \begin{figure}[t]
%   \centering 
%   \includegraphics[height=2.0in]{iccfdTemplate/exampleFigure.jpg}
%   \caption{This is the logo of ICCFD.}
%   \label{fig:logo}
% \end{figure}


% \subsection{Subsection Title Example}


% \subsubsection{Sub-subsection Title Example}


% \section{Conclusion and Future Work}
% ICCFD12 will be held at the Kobe international conference center in Port Island, Kobe, Japan. We look forward to welcoming you all to what we hope will be a memorable event.



% !TeX root = main.tex


\section{Introduction}
\label{sec:intro}

Computational fluid dynamics (CFD) simulation has been an efficient tool for investigation of flow problems, due to the fast development of computer hardware as well as numerical methods, such as finite difference (FD), finite volume (FV), finite element (FE) and spectral methods. Second-order FV method on unstructured grids \cite{van1979towards} is widely used in commercial and open-source CFD codes, thus being the workhorse of engineering flow simulations. Over recent decades, various high-order methods on unstructured grids have been developed, such as the high-order FV \cite{ollivier1997quasi_ENO,dumbser2007quadrature_WENO,wang2017compact_VR}, discontinuous Galerkin (DG) \cite{reed1973triangularDG,cockburn2001rungeDG}, hybrid FV/DG \cite{dumbser2008unified,li2022reconstructed,zhang2012class1}, residual distribution (RD)  \cite{abgrall2003construction} and flux reconstruction (FR) \cite{huynh2007flux,vincent2011new,wang2009unifying}. Such methods, compared with their second-order counterparts, have higher-order accuracy and lower dissipation/dispersion, while maintaining the capability of handling complex geometries. 
However, CFD solvers suffers from the deterioration of robustness when applied to challenging flow problems, even on high-quality computational meshes. A commonly encountered robustness issue in compressible flow simulation is the failure to preserve positivity of the density or pressure, which leads to nonphysical solutions or numerical instabilities.

Tremendous efforts have been made on developing positivity-preserving algorithms to improve robustness of numerical methods, especially high-order numerical methods, for compressible flow simulations. For the numerical methods based on explicit time integration, popular positivity-preserving algorithms can be categorized into two types \cite{xu2017bound}. The first type is the positivity-preserving scaling limiter \cite{zhang2010maximum,zhang2010positivity}, which compresses the solution polynomials to obtain positive density and pressure distributions, while being accuracy-preserving. It is proved that, by using the scaling limiter, a monotone numerical flux and a certain CFL condition, the FV and DG methods with strong stability preserving (SSP) Runge-Kutta time discretizations are positivity-preserving \cite{zhang2010positivity}. This approach has also been extended to FD methods \cite{zhang2012positivity,fan2022positivity}. The second type is the positivity-preserving flux limiter \cite{hu2013positivity,christlieb2015high,kuzmin2022bound}, which seeks a convex combination of the first-order monotone flux with the high-order flux that can achieve the positivity-preserving property under certain conditions. Compared with the scaling limiter, the advantage of the flux limiter is the easy application to high-order FD methods solving conservation laws and high-order schemes solving convection-diffusion problems \cite{xu2017bound}. However, the accuracy-preserving property of the flux limiters are often demonstrated by extensive numerical results \cite{xiong2013parametrized}, as it is difficult to be proved for general cases.  

Explicit time stepping methods suffer from the CFL constraint that restricts the time step sizes. 
For problems with very small allowable time step sizes, such as the high Reynolds number flows with large aspect ratio grids in the boundary layer region,  implicit time stepping methods are much more efficient.
Although most of the effort has been made for increasing accuracy of the time discretization and for increasing the efficiency of the nonlinear solver, only a few works exist in the literature concerning the positivity-preserving property of implicit methods \cite{qin2018implicit}, as it is difficult to verify if an implicit numerical scheme is positivity-preserving, even for a low-order one \cite{huang2024general}. This difficulty comes from the fact that, in implicit time stepping, the solutions are updated by solving a system of nonlinear equations iteratively. Parent \cite{parent2018positivity} developed a positivity-preserving dual-time stepping scheme for FV methods on structured grids solving Euler equations based on a Cauchy-Kowalevski procedure. Lian et al. \cite{lian2009solution} proposed a solution-limited time stepping scheme to enhance the reliability in steady-state compressible flow simulations. Qin et al. \cite{qin2018implicit} developed an implicit high-order positivity-preserving DG method for steady-state compressible Euler equations. Recently, Huang et al. \cite{huang2024general} proposed a positivity-preserving algorithm based on an iterative flux correction procedure, for finite volume methods simulating compressible flows on unstructured grids with a second-order backward difference (BDF2) dual-time stepping. 
%Based on the literature survey, it is found that it is of great interest to develop positivity-preserving algorithms for general high-order numerical methods on unstructured grids for unsteady compressible flow simulations.

In this paper, we develop a positivity-preserving algorithm for finite volume methods on unstructured grids with an implicit dual-time stepping, towards robust and efficient simulations of unsteady compressible flows.
In the positivity-preserving algorithm, admissible cell-averages are obtained by controlling physical and pseudo time step sizes which results in limited solution changes. To overcome the difficulty of unknown solution changes in time step limiting, we propose to use explicit time discretizations to obtain efficient estimations of future states. The allowable time step sizes are determined by limiting the relative solution changes to ensure positivity of updated solutions. Given positive cell-averages, admissible reconstruction polynomials can be obtained by applying a positivity-preserving scaling limiter \cite{zhang2010positivity}. The positivity-preserving algorithm is applied to a high-order variational finite volume method \cite{wang2017compact_VR} with an explicit first stage singly diagonally implicit Runge-Kutta (ESDIRK) time integration \cite{bijl2002implicitBDFvESDIRK}, and verified using a series of benchmark test cases. Numerical results demonstrate the robustness and high resolution of the positivity-preserving implicit high-order finite volume method. 

The remainder of this paper is organized as follows. Section \ref{sec:CFV} presents the high-order finite volume method on unstructured grids using implicit time stepping. Section \ref{sec:PP} presents the positivity-preserving algorithm for implicit finite volume methods. Numerical results are given in Section \ref{sec:Results} and concluding remarks are given in Section \ref{sec:Conclusions}.

\section{Implicit Finite Volume Method}
\label{sec:CFV}

\newcommand{\trans}{^\mathrm{T}}
\newcommand{\U}{\mathbf{U}}
\newcommand{\F}{\mathbf{F}}
\newcommand{\x}{\mathbf{x}}
\newcommand{\OO}{\mathbf{\Omega}}
\newcommand{\UM}{\overline{\U}}
\newcommand{\Fn}{\tilde{\F}}
\newcommand{\n}{\mathbf{n}}
\newcommand{\uu}{\overline{\mathbf{U}}}
\newcommand{\R}{\mathbf{R}}
\newcommand{\inc}{\mathrm\Delta}
\newcommand{\Tau}{\mathrm{T}}
\renewcommand{\real}{\mathrm{Re}}
\newcommand{\imag}{\mathrm{Im}}

\newcommand{\CFLt}{\text{CFL}_t}
\newcommand{\CFLtau}{\text{CFL}_\tau}

\newcommand{\eeqref}[1]{Eq.\eqref{#1}}
\newcommand{\um}{\overline{u}}
\newcommand{\us}{\mathbf{u}}
\newcommand{\SAll}{\mathcal{S}}

\newcommand{\FF}{\mathcal{F}}

\newcommand{\eye}{\mathbf{I}}

\subsection{Governing Equations}
\label{ssec:GovEq}

The Navier-Stokes equations governing unsteady compressible viscous flows can be expressed as
\begin{equation}
    \label{eq:NS}
    \dfrac{\partial \U}{\partial t} +
    \nabla \cdot (\F - \F_v)= 0,
\end{equation}
where $\U$ is the conservative variable vector, $\F$ is the inviscid flux tensor and $\F_v$
is the viscous flux tensor defined by
\begin{equation}
	\U = \begin{pmatrix}
		\rho \\ \rho \mathbf{u} \\ \rho E
	\end{pmatrix},\ \
	\F= \begin{pmatrix}
		\rho \mathbf{u}                \\
		\rho \mathbf{u}\otimes \mathbf{u}+p \mathbf{I}\\
		(\rho E+p)\mathbf{u}                  \\
	\end{pmatrix},\ \
	\F_{v} = \begin{pmatrix}
		0                                \\
		\boldsymbol{\tau}\\
	 	\boldsymbol{\tau} \cdot \mathbf{u} + \kappa \nabla T \\
	\end{pmatrix}.
\end{equation}
Here $\rho$ is the density, $\mathbf{u}$ is the velocity, $p$ is the pressure and $T$ is the temperature of the fluid. $\mathbf{I}$ is the identity tensor. $E$ is the total energy defined as
\begin{equation}
		E = \dfrac{1}{\gamma-1} \dfrac{p}{\rho} + \frac{1}{2} \mathbf{u} \cdot \mathbf{u},
\end{equation}
where $\gamma$ is the ratio of specific heat. The shear stress tensor $\boldsymbol{\tau}$ is defined as 
\begin{equation}
	\boldsymbol{\tau}= \mu \left[ \nabla \mathbf{u} + \left(\nabla \mathbf{u}\right)\trans - \frac{2}{3} \left(\nabla \cdot \mathbf{u}\right) \mathbf{I} \right],
\end{equation}
with $\mu$ being the dynamic viscosity. The heat conductivity is computed by $\kappa=C_p \mu /Pr$, where $C_p$ is the specific heat at constant pressure and $Pr$ is the Prandtl number. Equation \eqref{eq:NS} is closed by an equation of state
\begin{equation}
	p= \rho R T,
\end{equation}
where $R= \left(\gamma-1\right) C_p/\gamma$.
Neglecting viscous effects, i.e., $\mu=0$, the Navier-Stokes equations reduce to the Euler equations governing unsteady compressible inviscid flows.

\subsection{High-order Finite Volume Spatial Discretization}
\label{ssec:FV}

This subsection presents the general framework of high-order cell-centered finite volume method on unstructured grids for compressible flows.
The computational domain $\OO$ is partitioned into $N$ non-overlapping control volumes, i.e., $\OO= \cup^N_{i=1}{\OO_i}$. By integrating the governing equation \eqref{eq:NS} over control volume $\OO_i$, a semi-discrete finite volume scheme is obtained as
\begin{equation}
	\label{eq:Semi-FV}
	\dfrac{ \dd\UM_i }{\dd t} = -\dfrac{1}{\overline{\OO}_i} \oint_{\partial \OO_i} \left(\F - \F_v \right) \cdot \n \ \dd A,
\end{equation}
where $\overline{\OO}_i$ is the volume of $\OO_i$, $\partial \OO_i$ is the boundary of $\OO_i$ and $\n$ is the outward unit normal of $\partial \OO_i$, as shown in Figure \ref{fig:controlvolume}. The cell-average defined by
\begin{equation}
    \label{eq:FVMean}
    \UM_i \left(t\right)= \frac{1}{\overline{\OO}_i}\int_{\OO_i}\U \left(\x,t\right)\ \dd V,
\end{equation}
is the degree of freedom (DOF) of the finite volume method on cell $\OO_i$.
%where $\overline{\OO}_i$ is the volume of $\OO_i$.
The flux integral in \eqref{eq:Semi-FV} is computed by using a Gauss quadrature 
\begin{equation}
	\label{eq:Flux-Integral}
	\oint_{\partial \OO_i} \left(\F - \F_v \right) \cdot \n \ \dd A \approx \sum_{f \in \partial \OO_i} \sum_{g=1}^{N_g} w_g \left[\F \left(\U \left(\x_{f,g},t\right)\right) - \F_v \left(\U\left(\x_{f,g},t\right), \nabla \U \left(\x_{f,g},t\right)\right)\right] \cdot \n_f \ A_f,
\end{equation}
where $\n_f$ and $A_f$ are the outward unit normal and area of element interface $f$, respectively. $N_g$ is the number of quadrature points. $\x_{f,g}$ and $w_g$ are the position and weight of the $g$-th quadrature point on $f$, respectively. 

\begin{figure}[htbp!]
	\centering
	\includegraphics[width=0.25\linewidth]{pics/control_volume}
	\caption{A triangular element.}
	\label{fig:controlvolume}
\end{figure}


Given known cell-averages, a reconstruction is performed to obtain an approximate solution distribution on the computational domain $\OO$, to compute the states at quadrature points on element interfaces. Specifically, the solution on each control volume is approximated by a polynomial, i.e.,
\begin{equation}
    \label{eq:FVRec}
    \U_i(\x,t) = \UM_i \left(t\right) + \sum_{l=1}^{\mathrm{N_b}(k)}{\U_i^l \left(t\right) \varphi_{i,l}(\x) }, \quad \forall \ \x \in \OO_i, \quad i=1, \cdots, N,
\end{equation}
where $k$ is the degree of the polynomial, $\{\varphi_{i,l}(\x)\}$ are the polynomial basis functions and $\mathrm{N_b}(k)$ is the number of basis functions. The basis coefficients $\{\U^l_i\}$ are determined by using a reconstruction scheme, such as the variational reconstruction \cite{wang2017compact_VR} that will be presented in Section \ref{ssec:VR}. 

\begin{figure}[htbp!]
	\centering
	\includegraphics[width=0.3\linewidth]{pics/Riemann}
	\caption{Riemann problem on cell interface.}
	\label{fig:Riemann}
\end{figure}

Flux schemes based on solutions of Riemann problems are used to compute the numerical flux at the quadrature points in \eqref{eq:Flux-Integral}, since the piece-wise polynomial distribution is discontinuous across cell interfaces, as shown in Figure \ref{fig:Riemann}.  
The numerical flux is computed as
\begin{equation}
	\label{eq:num-flux}
	\left[\F \left(\U \left(\x_{f,g},t\right)\right) - \F_v \left(\U\left(\x_{f,g},t\right), \nabla \U \left(\x_{f,g},t\right)\right)\right] \cdot \n_f = \tilde{\F} \left(\U_L, \U_R,\n_f\right)- \tilde{\F}_v \left(\U_L, \nabla \U_L, \U_R, \nabla \U_R, \n_f\right),
\end{equation}
where $\tilde{\F}$ and $\tilde{\F}_v$ are the schemes used to compute the inviscid and viscous numerical fluxes, respectively, with the left and right states defined by
\begin{equation}
	\begin{dcases}
		&\U_L= \U_i\left(\x_{f,g},t \right), \quad \nabla \U_L= \nabla \U_i\left(\x_{f,g},t \right),\\
		&\U_R= \U_j\left(\x_{f,g} ,t \right), \quad \nabla \U_R= \nabla \U_j\left(\x_{f,g},t \right).
	\end{dcases}
\end{equation}
In this work, $\tilde{\F}$ is the local Lax-Friedrichs scheme and $\tilde{\F}_v$ is the viscous flux scheme by Wang et al. \cite{wang2017compact_VR}.

%In this work, the inviscid numerical flux is computed by using the local Lax-Friedrichs scheme
%\begin{equation}
%	\label{eq:LLF}
%	\F \left(\U \left(\x_{f,g},t\right)\right) \cdot \n_f = \dfrac{1}{2} \left[\F\left(\U_L\right) + \F\left(\U_R\right)\right]\cdot\n_f - \dfrac{1}{2} \max \left(\lambda\left(\U_L\right), \lambda\left(\U_R\right)\right) \left(\U_R - \U_L\right),
%\end{equation}
%where 
%\begin{equation}
%	\U_L= \U_i\left(\x_{f,g},t \right), \quad \U_R= \U_j\left(\x_{f,g} ,t \right), \quad \lambda\left(\U\right)= \abs{\mathbf{u} \cdot \n_f} + \sqrt{\gamma p/\rho},
%\end{equation}
%with $j$ being the face-neighboring cell that shares $f$ with cell $i$, i.e., $\OO_i \cap \OO_j= f$.
%The viscous numerical flux is computed by using the dGRP \cite{gassner2007contribution,gassner2008discontinuous} type scheme in \cite{wang2017compact_VR} 
%\begin{equation}
%		\label{eq:dGRP}
%	\F_v \left(\U \left(\x_{f,g},t\right), \nabla \U \left(\x_{f,g},t\right)\right) \cdot \n_f = \F_v \left(\tilde{\U}, \nabla \tilde{\U}  \right) \cdot \n_f,
%\end{equation}
%where 
%\begin{equation}
%	\tilde{\U} = \dfrac{1}{2} \left(\U_L+\U_R\right), \quad \nabla \tilde{\U}= \dfrac{1}{2}\left(\nabla \U_L+ \nabla \U_R\right) + \dfrac{1}{2 \Delta \tilde{x}} \left(\U_R - \U_L\right), \quad \Delta \tilde{x} = \dfrac{\min \left(\overline{\OO}_i,\overline{\OO}_j\right)}{A_f}.
%\end{equation}

The semi-discrete finite volume scheme \eqref{eq:Semi-FV} can be rewritten in an ordinary differential equation (ODE) form
\begin{equation}
	\label{eq:FVODECell}
	\derivative{\uu_i}{t} = \R_i \left(t, \left\{\U_j\right\}_{j \in S_i} \right),
\end{equation}
where can be integrated in time to update the cell-average in a step-by-step manner. The time integration method used in this work will be presented in Section \ref{ssec:TimeMarching}. 

\subsection{Variational Reconstruction}
\label{ssec:VR}

This subsection presents a variational reconstruction scheme \cite{wang2017compact_VR} that can achieve arbitrary high-order accuracy on a compact stencil involving only the current cell and its face-neighboring cells. The compact reconstruction stencil $S_i= \left\{i,j_1,j_2,j_3\right\}$ of a triangular element $i$ is shown in Figure \ref{fig:compactstencil}. In general, the conservative variables are reconstructed separately, allowing an implementation of solution reconstruction in a variable-by-variable manner. For the sake of presentation, the two-dimensional case is considered to illustrate the variational reconstruction procedure.

\begin{figure}[htbp!]
	\centering
	\includegraphics[width=0.3\linewidth]{pics/compact_stencil}
	\caption{Compact reconstruction stencil.}
	\label{fig:compactstencil}
\end{figure}

On control volume $\OO_i$, a conservative variable $u \in \U$ is approximated as
\begin{equation}
	\label{eq:recon_u}
	u_i \left(\x\right)= \overline{u}_i + \sum^{\mathrm{N_b}\left(k\right)}_{l=1} u_i^l \varphi_{i,l} \left(\x\right), 
\end{equation}
where $\{\varphi_{i,l}\}$ are the basis functions, $\{u^l_i\}$ are the unknown basis coefficients and $\mathrm{N_b}(k)$ is the number of basis functions. For a two-dimensional case, $\mathrm{N_b}(k)=  (k+2)(k+1)/2 -1$. In \eqref{eq:recon_u}, the time $t$ is omitted as it remains a constant during the reconstruction procedure.
%In order to determine the coefficients of polynomial bases $\U_i^l$ (or $u_i^l$ for each scalar) in
%\eeqref{eq:FVRec}, a reconstruction method needs to be specified.
%Traditional 2nd order FV methods for unstructured grid
%needs only to reconstruct a $k=1$ polynomial, namely linear
%distribution on each cell.
%The variational reconstruction \cite{wang2017compact_VR}
%is a compact high-order
%reconstruction scheme, and it achieves high-order
%accuracy on a compact stencil.
%The current section will explain the variational reconstruction briefly
%and specify details concerning following numerical tests.
In this work, the basis functions with the zero-mean property
\begin{equation}
	\frac{1}{\overline{\OO}_i} \int_{\OO_i} \varphi_{i,l} \left(\x\right) \ \dd V=0,
\end{equation}
are used to make the reconstruction polynomial \eqref{eq:recon_u} automatically satisfy the conservation condition
\begin{equation}
	\label{eq:zero-mean}
	\frac{1}{\overline{\OO}_i}\int_{\OO_i}u_i \left(\x\right)\ \dd V= \overline{u}_i.
\end{equation}
The zero-mean basis functions are defined by
\begin{equation}
        \varphi_{i,l} =
        \left(\frac{x - x_{i}}{\inc x_i}\right)^{p_l}
        \left(\frac{y - y_{i}}{\inc y_i}\right)^{q_l}
        -
        \overline{
            \left(\frac{x - x_{i}}{\inc x_i}\right)^{p_l}
            \left(\frac{y - y_{i}}{\inc y_i}\right)^{q_l}
        }, 
\end{equation}
with
\begin{equation}
	\overline{
		        \left(\frac{x - x_{i}}{\inc x_i}\right)^{p_l}
		        \left(\frac{y - y_{i}}{\inc y_i}\right)^{q_l}
		        }= \dfrac{1}{\overline{\OO}_i}\int_{\OO_i}   \left(\frac{x - x_{i}}{\inc x_i}\right)^{p_l}
		        \left(\frac{y - y_{i}}{\inc y_i}\right)^{q_l}  \dd V,
\end{equation} 
where $\left(x_i,y_i\right)$ and $\left(\inc x_i, \inc y_i\right)$ are the barycenter and characteristic length scales of cell $i$, respectively, and $\left(p_l,q_l\right)$ are the powers of the basis functions organized in ascending order of $p_l+q_l$. For instance, the powers of the cubic ($k=3$) reconstruction polynomial are 
\begin{equation}
	\left\{ \left(p_l,q_l\right) \right\}^9_{l=1}= \left\{ \left(1,0\right), \left(0,1\right), \left(2,0\right),\left(1,1\right),\left(0,2\right),\left(3,0\right),\left(2,1\right),\left(1,2\right),\left(0,3\right) \right\}.
\end{equation}
%The mean value term makes the basis zero-mean, which is calculated with:
%\begin{equation}
%    \overline{
%        \left(\frac{x - x_{c,i}}{\inc x_i}\right)^{p_l}
%        \left(\frac{y - y_{c,i}}{\inc y_i}\right)^{q_l}
%    }
%    =
%    \frac{1}{\overline{\OO}_j}\int_{\OO_j}{
%        \left(\frac{x - x_{c,i}}{\inc x_i}\right)^{p_l}
%        \left(\frac{y - y_{c,i}}{\inc y_i}\right)^{q_l}
%    }\dd \Omega
%\end{equation}
The characteristic length scales are used to non-dimensionalize the basis functions, to avoid a growth of the condition number of the reconstruction matrix during grid refinement \cite{abgrall1994essentially,friedrich1998weighted_WENO}. In this work, the length scales are set as $ \inc x_i = \inc y_i = \max_{\x \in \OO_i} \left\{\|\x - \x_{i}\|_2\right\}$. 
%\begin{equation}
%    \inc x_i = \inc y_i = \max_{\x \in \OO_i} \left\{\|\x - \x_{i}\|_2\right\}.
%\end{equation}

The objective of a reconstruction is to determine the unknown basis coefficients $u_i^l$,  $l=1,\cdots,\mathrm{N_b}\left(k\right)$, $i=1,\cdots,N$, given the cell-averages $\overline{u}_j$, $j=1,\cdots,N$. In the variational reconstruction \cite{wang2017compact_VR}, the linear equation system to determine the unknown basis coefficients is derived by minimizing a cost function using the variational method. Different cost function results in different reconstruction schemes. The cost function is defined as
\begin{equation}
	\label{eq:cost-function}
    I = \sum^{N_f}_{f=1} {I_f},
\end{equation}
where $I_f$ is an interfacial jump integration (IJI) on cell interface $f$ and $N_f$ is the total number of cell interfaces on the computational domain.
In this work, the IJI is defined as
\begin{equation}
	\label{eq:IJI}
    I_f = \omega_f^G \sum_{p+q=0}^{k} \int_{f}{
        \left[
            \omega_f^D(p,q)
            \left(
            \partialderivative{^{p+q}u_L}{x^p\partial y^q}
            -
            \partialderivative{^{p+q}u_R}{x^p\partial y^q}
            \right)
            \right]^2
        \dd A
    }
\end{equation}
where $L$ and $R$ are the two cells sharing interface $f$, $\omega_f^G$ is the geometric weight and
$\omega_f^D$ is the derivative weight. The IJI measures the jumps of the reconstruction polynomial and its spatial derivatives on cell interface. 
The derivative weights of the cubic variational reconstruction are
\begin{equation}
    \begin{aligned}
        \omega_f^D(0,0) & = \omega_D(0),                                                                                   \\
        \omega_f^D(1,0) & = \omega_f^D(0,1) = \omega_D(1) ,                                                                 \\
        \omega_f^D(2,0) & = \omega_f^D(0,2) = \omega_D(2),\ \omega_f^D(1,1) = \sqrt{2}\ \omega_D(2)  ,                 \\
        \omega_f^D(3,0) & = \omega_f^D(0,3) = \omega_D(3),\ \omega_f^D(1,2) = \omega_f^D(2,1) = \sqrt{3}\ \omega_D(3), \\
    \end{aligned}
    \label{eq:wdRotRatio}
\end{equation}
with
\begin{equation}
    \omega_D(0) = 1, \ \omega_D(1) = d_{LR}, \ \omega_D(2) = \frac{\left(d_{LR}\right)^2}{2}, \ \omega_D(3) = \frac{\left(d_{LR}\right)^3}{6},
    \label{eq:wdHQMOPT}
\end{equation}
where $d_{LR}= \left\| \x_L -\x_R \right\|_2$ is the distance between the barycenters of cell $L$ and $R$.
As only nearly uniform and isotropic meshes are used in the numerical simulations in this paper,
the geometric weights
$\omega^G_f$ are all set as 1.
The unknown basis coefficients are determined by minimizing the cost function \eqref{eq:cost-function}, resulting in a \textquote{smoothest} piece-wise polynomial distribution which has the smallest jumps measured by the IJIs on cell interfaces. 
Linear reconstruction equations can be derived by using the variational method, i.e.,
\begin{equation}
	\label{eq:minimization}
    \partialderivative{I}{u_i^l} = 0, \ l=1,\cdots, \mathrm{N_b}(k),\  i=1,\cdots N.
\end{equation}
By substituting \eqref{eq:recon_u} into \eqref{eq:IJI}, \eqref{eq:cost-function} and \eqref{eq:minimization}, a linear equation system is obtained as following
\begin{equation}
    \label{eq:vrBlockEq}
  \mathbf{A}_{i} \us_i
    =
    \sum_{j \in S_i, j \neq i} \mathbf{B}^j_{i} \us_j + \mathbf{b}_{i}, \quad i=1, \cdots, N,
\end{equation}
where the elements of matrices $\mathbf{A}_i$, $\mathbf{B}^j_i \in \mathbb{R}^{\mathrm{N_b}(k) \times \mathrm{N_b}(k)}$ and vectors $\mathbf{u}_i$, $\mathbf{u}_j$, $\mathbf{b}_i \in \mathbb{R}^{\mathrm{N_b}(k)}$ are
\begin{equation}
    \label{eq:vrCoeffs}
    \begin{aligned}
        &\mathbf{A}_{i} [m,n]=
               \sum_{j \in S_i, j \neq i} \omega_f^G\int_{\OO_i \cap \OO_j}{
                    \sum_{p+q=0}^{k}
                    \left(\omega_f^D(p,q)\right)^2
                    \partialderivative{^{p+q}\varphi_{i,m}}{x^p\partial y^q}
                    \partialderivative{^{p+q}\varphi_{i,n}}{x^p\partial y^q}
                    \ \dd A},
           \\
        &\mathbf{B}^j_{i} [m,n]= 
                \omega_f^G\int_{\OO_i \cap \OO_j}{
                    \sum_{p+q=0}^{k}
                    \left(\omega_f^D(p,q)\right)^2
                    \partialderivative{^{p+q}\varphi_{i,m}}{x^p\partial y^q}
                    \partialderivative{^{p+q}\varphi_{j,n}}{x^p\partial y^q}
                    \ \dd A},
           \\
        &\mathbf{b}_{i} [m]=
        			\sum_{j \in S_i, j \neq i} 
                    \omega_f^G \left(\omega_f^D(0,0)\right)^2
                    \int_{\OO_i \cap \OO_j}{
                        \varphi_{i,m} \left(\overline{u}_j - \overline{u}_i\right)
                        \ \dd A},
		 \\
		 & \mathbf{u}_i [m]= u^m_i, \\ 
		 & \mathbf{u}_j [m]= u^m_j.
    \end{aligned}
\end{equation}
It is observed from \eqref{eq:vrBlockEq} that the variational reconstruction is implicit, as the unknown coefficients of the face-neighboring cells are required to determine the unknown coefficients of the current cell. 
By assembling the linear equations \eqref{eq:vrBlockEq} of all cells, we obtain a global linear equation system
\begin{equation}
	\label{eq:global-system}
	\mathbf{A} \mathbf{u} = \mathbf{b},
\end{equation}
where 
\begin{equation}
	\begin{aligned}
		&\mathbf{A}= \mathbf{D} -  \mathbf{L} -  \mathbf{U}, 
		\ \mathbf{D}= \left\{\mathbf{A}_i\right\}, 
		\ \mathbf{L}= \left\{\mathbf{B}^j_i, j<i\right\},
		\ \mathbf{U}= \left\{\mathbf{B}^j_i, j>i\right\},
		\\
		& \mathbf{u}= \left\{\mathbf{u}_i\right\}, \ \mathbf{b}= \left\{\mathbf{b}_i\right\}.
	\end{aligned}
\end{equation}
It is proved in \cite{wang2017compact_VR} that the large and sparse reconstruction matrix $\mathbf{A}$ is symmetric and positive definite, which guarantees the existence and uniqueness of the solution of the linear equation system \eqref{eq:global-system}. This is a significant advantage of the variational reconstruction over other existing high-order reconstructions on unstructured grids.
The linear equation system \eqref{eq:global-system} is solved iteratively using the block Gauss-Seidel method, of which the convergence is proved in \cite{wang2017compact_VR}.
Each block Gauss-Seidel iteration is compact as it only relies on the information of the current and face-neighboring cells.
%It is proved in \cite{wang2017compact_VR} that the block Gauss-Seidel method is convergent for the linear equation system \eqref{eq:global-system}.

In the reconstructions of solutions involving discontinuities, a WBAP limiter \cite{li2011multi,li2012multi} is used to suppress spurious oscillations. The limiting process on the element $i$ in Figure \ref{fig:compactstencil} is presented for illustration. The limited reconstruction polynomial
\begin{equation}\label{eqn:Limited_Polynomial}
	\tilde{u}_{i} \left(\x\right) = \overline{u}_{i} + \sum_{l=1}^{\mathrm{N_b}\left(k\right)} \tilde{u}_i^l \varphi_{i,l}  \left(\x\right)
\end{equation}
is a nonlinear weighted average of the reconstruction polynomial $u_i \left(\x\right)$ and three additional candidate polynomials $u_{j \rightarrow i} \left(\x\right), j=j_1, j_2, j_3$ obtained through a secondary reconstruction \cite{li2012high}. The limited coefficients $\tilde{u}^l_i$ is computed by
\begin{equation}\label{eq:Limiting}
	\tilde{u}_i^l= L\left(u_i^l, u^l_{j_1 \rightarrow i}, u^l_{j_2 \rightarrow i}, u^l_{j_3 \rightarrow i}\right),
\end{equation} 
with a limiting function
\begin{equation}
	L(a_0,a_1,...,a_J)=a_0\cdot{W\left(1,\frac{a_1}{a_0},...,\frac{a_J}{a_0}\right)}, \quad W=\frac{n+\sum_{m=1}^{J}{1/\theta_m^{p-1}}}{n+\sum_{m=1}^{J}{1/\theta_m^{p}}},
\end{equation}
where $n=10$ and $p=4$. The limiting is performed in a successive manner in characteristic space \cite{li2012multi}.

\subsection{Implicit Time Marching}
\label{ssec:TimeMarching}

As mentioned in the previous subsection, the linear equation system of the variational reconstruction is solved iteratively. The reconstruction will be very expensive if the iteration needs to reach convergence at each single time step. A reconstruction and implicit dual-time stepping coupled iteration procedure has been proposed in \cite{wang2016compact1_VR} to address this computational efficiency issue. In the coupled iteration procedure, the reconstruction iteration is performed only once at each pseudo time step. The coupling between the reconstruction and the time integration makes these two procedures achieve convergence synchronously. By using the coupled iteration, the implicit nature of the variational reconstruction does not result in additional cost, thus ensuring the high computational efficiency of the variational finite volume method. 

In this work, the time integration scheme for unsteady flow simulation is the six-stage, fourth-order explicit first stage singly diagonal implicit Runge-Kutta (ESDIRK4) \cite{bijl2002implicitBDFvESDIRK}. By using this method, \eqref{eq:FVODECell} is integrated in time as
\begin{subequations}\label{eq:ESDIRK4}
		 \begin{align}
			&\uu^{n,1}_i \ = \uu^{n}_i, \label{eq:esdirk4-1}\\        
			&\uu^{n,s}_i \ = \uu^{n}_i + \inc t^n \sum_{q=1}^{s} a_{sq} \R_i^{n,q},\ \ s = 2, \cdots, 6, \label{eq:esdirk4-s} \\
			&\uu^{n+1}_i = \uu^{n,6}_i, \label{eq:esdirk4-6}
		\end{align}
\end{subequations}
%\begin{equation}
%    \begin{aligned}
%    	&\uu^{n,1}_i \ = \uu^{n}_i, \\        
%        &\uu^{n,s}_i \ = \uu^{n}_i + \inc t^n \sum_{q=1}^{s} a_{sq} \R_i^{n,q},\ \ s = 2, \cdots, 6 \\
%        &\uu^{n+1}_i = \uu^{n,6}_i, \\
%    \end{aligned}
%\end{equation}
where
\begin{equation}
	\R^{n,q}_i  = \R_i \left(t^n + c_q \inc t^n, \left\{\U^{n,q}_j\right\}_{j \in S_i}\right),
\end{equation}
with $\inc t^n= t^{n+1} - t^n$ being the time step size, which is in general a user-specified constant.
The Runge-Kutta coefficients $a_{sq}$ and $c_q$ of ESDIRK4 can be found
in \cite{bijl2002implicitBDFvESDIRK}.

The implicit and nonlinear equation \eqref{eq:esdirk4-s} is solved iteratively using a dual-time stepping technique, in which a pseudo-time variable $\tau$ is introduced. The pseudo-time integration scheme for  \eqref{eq:esdirk4-s} is
\begin{equation}
    \dfrac{\uu^{n,s,m+1}_i-\uu^{n,s,m}_i}{\inc \tau_i} + \dfrac{\uu^{n,s,m+1}_i - \uu^{n}_i}{\inc t^n} = \sum_{q=1}^{s-1} a_{sq} \R^{n,q}_i + a_{ss} \R_i^{n,s,m+1},
\end{equation}
where $m$ is the index for pseudo-time step, and
\begin{equation}
	\R^{n,s,m+1}_i  = \R_i \left(t^n + c_s \inc t^n, \left\{\U^{n,s,m+1}_j\right\}_{j \in S_i}\right).
\end{equation}
By using the linear approximation
\begin{equation}
	\R^{n,s,m+1}_i \approx \R^{n,s,m}_i + \sum_{j \in S_i} \dfrac{\partial \R_i}{\partial \UM_j} \left(\uu^{n,s,m+1}_j-\uu^{n,s,m}_j\right),
\end{equation}
we obtain a linear equation system to update the cell-average in pseudo-time direction as following
\begin{equation}
%	\label{eq:pseudo-time-equation}
    \left(\frac{\eye}{\inc \tau_i} + \frac{\eye}{\inc t^n} -a_{ss}\partialderivative{\R_i }{\uu_i} \right) \inc \uu^{n,s, m}_i
    -
    a_{ss}\sum_{j\in S_i,j\neq i} {
        \partialderivative{\R_i}{\uu_j} \inc \uu^{n,s, m}_j
    }
    = \tilde{\R}^{n,s,m}_i,
    \label{eq:linearTauUpdate}
\end{equation}
where
\begin{equation}
	\label{eq:define-inc}
    \inc \uu^{n,s,m}_i= \uu^{n,s,m + 1}_i - \uu^{n,s,m}_i, \quad \tilde{\R}^{n,s,m}_i= \sum_{q=1}^{s-1} a_{sq} \R^{n,q}_i + a_{ss} \R_i^{n,s,m} - \dfrac{\uu^{n,s,m}_i - \uu^{n}_i}{\inc t^n}.
\end{equation}
The linear equation system \eqref{eq:linearTauUpdate} is solved by using the matrix-free LU-SGS approach \cite{luo1998fast}. The solution is advanced in pseudo-time until $\uu^{n,s,m+1}_{i}$ converges to $\uu^{n,s}_{i}$, forming an inner iteration at the $s$-th stage of ESDIRK4. The convergence criterion of the inner iteration is that the $L_1$ norm of the pseudo-time derivative decreases by a certain number of orders of magnitude.   

The local pseudo-time step size is computed by
\begin{equation}
	\inc \tau_{i} = \frac{\CFLtau \overline{\OO}_i }
	{\sum_{f \in \partial \OO_i}{A_f}\lambda_{f}}
\end{equation}
where $\lambda_{f}$ is the spectral radius estimated on cell interface $f$ and $\CFLtau$ is the CFL number used to control the convergence speed of inner iteration.

%The pseudo time step is determined locally on each cell to account for
%the different time scales in the computational region.
%The current paper uses first order LLF approximate evaluation
%of the Jacobian terms in \eeqref{eq:linearTauUpdate} and uses matrix-free
%LU-SGS to solve the increment values \cite{luo1998fast,luo2001accurate},
%which grants enough convergence in unsteady problems. 
%The pseudo time iteration is stopped when
%\begin{equation}
%    \|\FF^{n,s,m}\| \leq \varE_{\tau} \|\FF^{n,s,0}\|
%\end{equation}
%and $\FF^{n,s,m}$ is the globally assembled residual vector from 
%$\FF^{n,s,m}_i$, $\varE_{\tau}$ is the convergence threshold for
%the nonlinear residual $\FF^{n,s,m}$, and $\|\cdot\|$ is 
%a norm. For transient flow the current research finds setting $\varE_{\tau}$
%from $10^{-3}$ to $10^{-4}$ sufficiently accurate.

\section{Positivity-Preserving Algorithm}
\label{sec:PP}

In compressible flow simulations, density and pressure need to be preserved positive, to avoid nonphysical solutions or numerical instabilities. Therefore, a positivity-preserving property is desired for numerical methods, especially high-order numerical methods, to perform robust simulations of compressible flows.

The set of admissible states is defined by
\begin{equation}
	G= \left\{\U =  \begin{pmatrix}
		\rho \\ \rho \mathbf{u} \\ \rho E
	\end{pmatrix} \middle| \rho>0 \ \mathrm{and} \ p= 
	\left(\gamma-1\right) \left(\rho E - \dfrac{1}{2} \rho \mathbf{u} \cdot \mathbf{u}\right) >0 \right\}.
\end{equation}
The density $\rho$ is a linear function of $\U$, and satisfies
\begin{equation}
	\label{eq:density-linear}
		\rho\left(w\U_1+(1-w)\U_2\right) = w\rho\left(\U_1\right) + (1-w) \rho \left(\U_2\right) , \quad 0 \leq w \leq 1.
\end{equation}
The pressure $p$ is a concave function of $\U$ if $\rho\left(\U\right) \geq 0$, and satisfies Jensen's inequality
\begin{equation}
	\label{eq:pressure-Jensen}
		p \left(w\U_1+(1-w)\U_2\right)  \geq w p\left(\U_1\right) + (1-w) p \left(\U_2\right), \quad 0 \leq w \leq 1, \quad \mathrm{if} \ \rho\left(\U_1\right) \geq 0, \ \rho\left(\U_2\right) \geq 0.
\end{equation}
A positivity-preserving numerical scheme satisfies that $\U^{n+1} \in G$, if $\U^{n} \in G$.  In a finite volume method, the numerical solution at $t^{n+1}$ is
\begin{equation}
%	\label{eq:FVRec}
	\U^{n+1}_i (\x)= \UM^{n+1}_i + \sum_{l=1}^{\mathrm{N_b}(k)}{\U_i^{n+1,l} \varphi_{i,l}(\x) }, \quad i=1, \cdots,N.
\end{equation}
An implicit finite volume solution $\U^{n+1} \in G$ can be obtained through the following two steps:
\begin{enumerate}[label=(\arabic*)]
	\item Given $\U^n \in G$, obtain $\UM^{n+1} \in G$ by using a proper implicit time-stepping scheme;
	\item Given $\UM^{n+1} \in G$, obtain $\U^{n+1} \in G$ by using a proper limiting algorithm.
\end{enumerate}  
The second step can be accomplished by using the scaling limiter of Zhang et al. \cite{zhang2010positivity}. Therefore, the key of developing a positivity-preserving implicit finite volume method is to design an implicit time-stepping scheme that leads to admissible updated cell-averages.

In this paper, we propose a positivity-preserving algorithm for finite volume methods with implicit dual-time discretizations. In this algorithm, the physical and pseudo time step sizes are controlled by limiting the solution changes, to obtain admissible cell-averages at the next physical and pseudo time steps, respectively. The details of the proposed positivity-preserving algorithm will be presented in the following subsections.

%The finite volume method with implicit time marching
%is able to steadily produce numerical results in
%a wide range of flow problems.
%However, when flow conditions are extreme,
%the numerical process could fail due to
%emergence of invalid states including
%negative density and negative internal energy.
%High-order finite volume methods suffer from this issue especially,
%which could even fail due to the initial transients
%of a low Mach flow problem.
%This kind of numerical failure is referred to
%as positivity-preserving problems in the current research.
%
%The current section will introduce a series of simple
%and straightforward algorithms
%to mitigate the change of encountering positivity-preserving problems.

\subsection{Physical Time Step Limiting}
\label{ssec:physical-limiting}

In an unsteady flow simulation using implicit time marching, the physical time step size is usually set as a user-specified constant, i.e., $\inc t^n= \inc t_{max}$, based on a prior estimation of the time scale of the problem. For challenging problems, there are transient states with very low density or pressure, making it difficult for the numerical method to preserve positivity using the original time step size $\inc t_{max}$.
The positivity-preserving difficulty may be overcome by reducing the time step size temporarily. Therefore, we propose to control the physical time step size dynamically based on a positivity-preserving principle
\begin{equation}
		\label{eq:ppLimit} 
%		\begin{aligned}
			\rho\left(\uu_i^{n+1} \right)>0, \quad 
			p\left(\uu_i^{n+1}\right)>0 ,
%		\end{aligned}
\end{equation}
for all cells $i=1, \cdots,N$.
However, the solution $\UM_i^{n+1}$ is unknown at the beginning of the time step. To overcome this difficulty, we propose an solution estimation based on a forward Euler time stepping, which is given by
\begin{equation}
	\label{eq:estimation-dt}
	\UM_i^{n+1} \approx  \uu_i^n + \inc t^{n}  \R_i^n.
\end{equation}
 A time step size $ \inc t^{n}_{pp}$ satisfying the positivity-preserving principle \eqref{eq:ppLimit} can be determined by limiting the solution changes as follows
\begin{equation}
    \label{eq:ppLimitDt}
    \begin{aligned}
        \rho\left(\uu_i^n + \inc t^{n}_{pp}\R_i^n\right)
         & \geq
        (1- \eta_t)
        \rho\left(\uu_i^n\right)>0, \\
        p\left(\uu_i^n + \inc t^{n}_{pp}\R_i^n\right)
         & \geq
        (1- \eta_t)
        p\left(\uu_i^n\right)>0,
    \end{aligned}
\end{equation}
for all cells $i=1, \cdots,N$, where $\eta_t\in(0,1)$ is the parameter used to control the relative solution changes. It is noted that the forward Euler scheme is first-order accurate, and thus the error of the estimation in \eqref{eq:estimation-dt} is $\mathcal{O}(\inc t^2)$. As the allowable solution changes in \eqref{eq:ppLimitDt} are $\eta_t \rho\left(\uu_i^n\right) \sim \mathcal{O}(1)$ and $\eta_t p\left(\uu_i^n\right) \sim \mathcal{O}(1)$, the positivity of $\UM_i^{n+1}$  can be preserved, provided that $\eta_t$ is not overly large.

The allowable time step size for cell $i$ can be determined in two steps according to \eqref{eq:ppLimitDt}. First, we assume an allowable time step size $\alpha_{t,i,\rho} \inc t_{max}$, with $\alpha_{t,i,\rho} \in (0,1]$, to preserve positivity of density. According to \eqref{eq:density-linear} and \eqref{eq:ppLimitDt}, we have
\begin{equation}
	\rho\left(\UM^n_i + \alpha_{t,i,\rho} \inc t_{max} \R^n_i\right)= \left(1- \alpha_{t,i,\rho}\right)\rho\left(\UM^n_i\right) + \alpha_{t,i,\rho} \rho\left(\UM^n_i + \inc t_{max} \R^n_i\right)  \geq
	(1- \eta_t)
	\rho\left(\uu_i^n\right),
\end{equation}
and thus obtain
\begin{equation}
	\alpha_{t,i,\rho} = \begin{dcases}
		 1, & \ \mathrm{if} \ \rho\left(\UM^n_i + \inc t_{max} \R^n_i\right) \geq  \rho\left(\UM^n_i\right), \\
		\min\left(1, \dfrac{-\eta_t \rho\left(\UM^n_i\right)}{\rho \left(\UM^n_i + \inc t_{max} \R^n_i\right)-  \rho\left(\UM^n_i\right)} \right), &\ 
		%\mathrm{if} \ \rho\left(\UM^n_i + \inc t_{max} \R^n_i\right)<  \rho\left(\UM^n_i\right).
		\mathrm{else}.
	\end{dcases}
\end{equation}
Second, we assume an allowable time step size $\alpha_{t,i,p}  \alpha_{t,i,\rho} \inc t_{max}$, with $\alpha_{t,i,p} \in (0,1]$, to preserve positivity of pressure. According to \eqref{eq:pressure-Jensen} and \eqref{eq:ppLimitDt}, we have
\begin{equation}
	p\left(\UM^n_i +  \alpha_{t,i,p} \alpha_{t,i,\rho} \inc t_{max} \R^n_i\right) \geq \left(1-\alpha_{t,i,p}\right)p\left(\UM^n_i\right) + \alpha_{t,i,p} p\left(\UM^n_i + \alpha_{t,i,\rho} \inc t_{max} \R^n_i\right) \geq
	(1- \eta_t)
	p\left(\uu_i^n\right),
\end{equation}
and thus obtain a sufficient condition
\begin{equation}
	\alpha_{t,i,p} = \begin{dcases}
		1, & \ \mathrm{if} \ p \left(\UM^n_i + \alpha_{t,i,\rho} \inc t_{max} \R^n_i\right) \geq  p \left(\UM^n_i\right), \\
		\min\left(1, \dfrac{-\eta_t p \left(\UM^n_i\right)}{p \left(\UM^n_i +  \alpha_{t,i,\rho} \inc t_{max} \R^n_i\right)-  p \left(\UM^n_i\right)} \right), &\ %\mathrm{if} \ p \left(\UM^n_i + \alpha_{t,i,\rho} \inc t_{max} \R^n_i\right)<  p\left(\UM^n_i\right).
		\mathrm{else}.
	\end{dcases}
\end{equation}
The global time step size is computed by
\begin{equation}
    \inc t^n = \alpha_t \inc t_{max} 
\end{equation}
where
\begin{equation}
        \alpha_t = \min_i(\alpha_{t,i,\rho}\alpha_{t,i,p}).
\end{equation}
It can be observed from the above equations that $\inc t^n \leq \inc t_{max}$, as $0<\alpha_t \leq 1$.

\subsection{Pseudo Time Step Limiting}
\label{ssec:pseudo-limiting}

The physical time step limiting strategy described in the previous subsection is used to obtain a solution $\UM_i^{n+1} \in G$, which is the converged solution of the last stage of ESDIRK4 \eqref{eq:ESDIRK4}. However, intermediate solutions with negative density or pressure might be encountered during the inner iteration at each stage. Therefore, we propose to limit the solution changes by controlling the pseudo-time step sizes during the inner iteration process, similar to the limiting strategy in physical time direction in Section \ref{ssec:physical-limiting}. 

For the $m$-th inner iteration at the $s$-stage of ESDIRK4, we have the following positivity-preserving requirement
\begin{equation}
	\label{eq:ppLimit-Tau} 
	%		\begin{aligned}
		\rho\left(\uu_i^{n,s,m+1} \right)>0, \quad 
		p\left(\uu_i^{n,s,m+1}\right)>0,
		%		\end{aligned}
\end{equation}
for all cells $i=1,\cdots,N$. However, $\UM^{n,sm+1}_i$ is unknown at the beginning of the pseudo-time step. Based on an explicit dual-time  discretization of \eqref{eq:esdirk4-s} given by
\begin{equation}
	\dfrac{\uu^{n,s,m+1}_i-\uu^{n,s,m}_i}{\inc \tau_i} + \dfrac{\uu^{n,s,m}_i - \uu^{n}_i}{\inc t^n} \approx \sum_{q=1}^{s-1} a_{sq} \R^{n,q}_i + a_{ss} \R_i^{n,s,m},
\end{equation}
we obtain an solution estimation
\begin{equation}
	\label{eq:estimation-dtau}
	\UM_i^{n,s,m+1} \approx  \UM_i^{n,s,m} + \inc \tau_i  \tilde{\R}_i^{n,s,m},
\end{equation}
with $ \tilde{\R}_i^{n,s,m}$ defined in \eqref{eq:define-inc}.
A pseudo-time step size $ \inc \tau_{i,pp}$ satisfying the positivity-preserving principle \eqref{eq:ppLimit-Tau} can be determined by limiting the solution changes as follows
\begin{equation}
	\label{eq:ppLimitDtau}
	\begin{aligned}
		\rho\left(\uu_i^{n,s,m} +  \inc \tau_{i,pp} \tilde{\R}_i^{n,s,m}\right)
		& \geq
		(1- \eta_\tau)
		\rho\left(\uu_i^{n,s.m}\right)>0, \\
		p\left(\uu_i^{n,s,m} + \inc  \tau_{i,pp} \tilde{\R}_i^{n,s,m}\right)
		& \geq
		(1- \eta_\tau)
		p\left(\uu_i^{n,s,m}\right)>0,
	\end{aligned}
\end{equation}
where $\eta_\tau\in(0,1)$ is the parameter used to control the relative solution changes.

The allowable pseudo-time step size for cell $i$ can be determined according to \eqref{eq:ppLimitDtau} using the same two-step procedure in Section \ref{ssec:physical-limiting}. The allowable time step size is $\inc \tau_{i,pp}= \alpha_{\tau,i,p}\alpha_{\tau,i,\rho} \inc  \tau_{i}$, 
%with $\alpha_{\tau,i,\rho}$, $\alpha_{\tau,i,p} \in (0,1]$, 
where
%\begin{equation}
%	\rho\left(\UM^n_i + \alpha_{t,i,\rho} \inc t_{max} \R^n_i\right)= \left(1- \alpha_{t,i,\rho}\right)\rho\left(\UM^n_i\right) + \alpha_{t,i,\rho} \rho\left(\UM^n_i + \inc t_{max} \R^n_i\right)  \geq
%	(1- \eta_t)
%	\rho\left(\uu_i^n\right),
%\end{equation}
%and thus obtain
\begin{equation}
	\alpha_{\tau,i,\rho} = \begin{dcases}
		1, & \mathrm{if} \ \rho\left(\UM^{n,s,m}_i + \inc \tau_{i} \tilde{\R}^{n,s,m}_i\right) \geq  \rho\left(\UM^{n,s,m}_i\right), \\
		\min\left(1, \dfrac{-\eta_\tau \rho\left(\UM^{n,s,m}_i\right)}{\rho \left(\UM^{n,s,m}_i + \inc \tau_{i} \tilde{\R}^{n,s,m}_i\right)-  \rho\left(\UM^{n,s,m}_i\right)} \right), &
		%\mathrm{if} \ \rho\left(\UM^{n,s,m}_i + \inc \tau_{i} \tilde{\R}^{n,s,m}_i\right) <  \rho\left(\UM^{n,s,m}_i\right).
		\mathrm{else},
	\end{dcases}
\end{equation}
%Second, we assume an allowable time step size $\alpha_{t,i,p}  \alpha_{t,i,\rho} \inc t_{max}$, with $\alpha_{t,i,p} \in (0,1]$, to preserve positivity of pressure. According to \eqref{eq:pressure-Jensen} and \eqref{eq:ppLimitDt}, we have
%\begin{equation}
%	p\left(\UM^n_i +  \alpha_{t,i,p} \alpha_{t,i,\rho} \inc t_{max} \R^n_i\right) \geq \left(1-\alpha_{t,i,p}\right)p\left(\UM^n_i\right) + \alpha_{t,i,p} p\left(\UM^n_i + \alpha_{t,i,\rho} \inc t_{max} \R^n_i\right) \geq
%	(1- \eta_t)
%	p\left(\uu_i^n\right),
%\end{equation}
%and thus obtain
and 
\begin{equation}
	\alpha_{\tau,i,p} = \begin{dcases}
		1, & \mathrm{if} \ p \left(\UM^{n,s,m}_i + \alpha_{\tau,i,\rho} \inc \tau_{i} \tilde{\R}^{n,s,m}_i\right) \geq  p \left(\UM^{n,s,m}_i\right), \\
		\min\left(1, \dfrac{-\eta_\tau p \left(\UM^{n,s,m}_i\right)}{p \left(\UM^{n,s,m}_i +  \alpha_{\tau,i,\rho} \inc \tau_{i} \tilde{\R}^{n,s,m}_i\right)-  p \left(\UM^{n,s,m}_i\right)} \right), &
		%\mathrm{if} \ p \left(\UM^{n,s,m}_i + \alpha_{\tau,i,\rho} \inc \tau_{i} \tilde{\R}^{n,s,m}_i\right)<  p\left(\UM^{n,s,m}_i\right).
		\mathrm{else}.
	\end{dcases}
\end{equation}

%The size of pseudo time step is determined
%by the CFL condition first:
%\begin{equation}
%    \inc \tau_{i,CFL} = \frac{\CFLtau \overline{\OO_i} }
%    {\sum_{j\in S_i, j\neq i}{\overline{f_{ij}}\lambda_{ij}}}
%\end{equation}
%where $\CFLtau$ is the CFL number used to control
%pseudo time stepping, $\overline{\OO_i}$ and $\overline{f_{ij}}$ 
%are volume and area of cell and interfaces respectively, 
%$\lambda_{ij}$ is the spectral radius estimated on each interface.
%By controlling $\CFLtau$, the convergence speed and 
%stability of iteration can be improved. 
%In positivity-preserving troubled cases, 
%such control on $\inc \tau_i$ is insufficient.
%The current research proposes to use a limiting 
%process similar to that of physical time step 
%to decide the pseudo time step used for 
%iteration $m$:
%\begin{equation}
%    \label{eq:ppLimitDtau}
%    \begin{aligned}
%        \rho\left(\uu_i^{n,s,m} + \inc \tau_{i, pp}\FF_i^{n,s,m}\right)
%         & \geq
%        (1- \eta_\tau)
%        \rho\left(\uu_i^{n,s,m}\right) \\
%        e\left(\uu_i^{n,s,m} + \inc \tau_{i, pp}\FF_i^{n,s,m}\right)
%         & \geq
%        (1- \eta_\tau)
%        e\left(\uu_i^{n,s,m}\right), \ \ \forall i
%    \end{aligned}
%\end{equation}
%Similar to $\eta_t$, $\eta_\tau\in(0,1)$ is a relaxation factor. 
%Once again, a sufficient condition is
%\begin{equation}
%    \inc\tau_{i,pp} =  \alpha_{\tau,i} \inc\tau_{i,CFL}
%\end{equation}
%where
%\begin{equation}
%    \alpha_{\tau,i} = \min\left(
%        \eta_\tau\frac{\rho(\uu_i^{n,s,m})}{-\rho_-(\inc\tau_{i,CFL} \FF_i^{n,s,m})},
%        \eta_\tau\frac{e(\uu_i^{n,s,m})}{-e_-(\inc\tau_{i,CFL} \FF_i^{n,s,m})},
%        1
%        \right)
%\end{equation}
%Note that the reference time scale here is $\inc\tau_{i,CFL}$, 
%and the compressing of $\inc\tau_i$ is non-uniform over the cells.

It is noted that $\UM^{n,s,m+1}_i$ is obtained by performing one LU-SGS iteration for \eqref{eq:linearTauUpdate}, thus being not the converged solution of the linear equation system. Therefore, there is not theoretical result for the error of estimation in \eqref{eq:estimation-dtau}, and thus not guarantee for the positivity of the solution $\UM^{n,s,m+1}_i$. To address this issue, an increment correction will be introduced in the next subsection.

\subsection{Increment Correction}

As pointed out in the previous subsection, although the physical and pseudo time steps are limited, the intermediate results may have negative density or pressure before the convergence of the LU-SGS iteration.
Therefore, a pseudo-time increment correction is designed to ensure the positivity of the updated solutions.
The increment correction is based on the following conditions
\begin{equation}
    \label{eq:ppLimitInc}
    \begin{aligned}
        \rho\left(\uu_i^{n,s,m} + \alpha_{\inc,i} \inc \uu_i^{n,s,m}\right)
         & \geq
        (1- \eta_\inc)
        \rho\left(\uu_i^{n,s,m}\right), \\
        p\left(\uu_i^{n,s,m} + \alpha_{\inc,i} \inc \uu_i^{n,s,m}\right)
         & \geq
        (1- \eta_\inc)
        p\left(\uu_i^{n,s,m}\right),
    \end{aligned}
\end{equation}
where $\eta_\inc \in (0,1)$ is a relaxation parameter to control the relative solution changes.
By using the same two-step approach used in physical/pseudo time step limiting, the relaxation parameter is computed by
\begin{equation}
    \alpha_{\inc,i}= \alpha_{\inc,i,\rho} \alpha_{\inc,i,p},
\end{equation} 
where
\begin{equation}
	\alpha_{\inc,i,\rho} = \begin{dcases}
		1, & \mathrm{if} \ \rho\left(\UM^{n,s,m}_i +\inc \uu_i^{n,s,m}\right) \geq  \rho\left(\UM^{n,s,m}_i\right), \\
		\min\left(1, \dfrac{-\eta_\inc \rho\left(\UM^{n,s,m}_i\right)}{\rho \left(\UM^{n,s,m}_i +\inc \uu_i^{n,s,m} \right)-  \rho\left(\UM^{n,s,m}_i\right)} \right), &
		%\mathrm{if} \ \rho\left(\UM^{n,s,m}_i + \inc \uu_i^{n,s,m}\right) <  \rho\left(\UM^{n,s,m}_i\right).
		\mathrm{else},
	\end{dcases}
\end{equation}
%Second, we assume an allowable time step size $\alpha_{t,i,p}  \alpha_{t,i,\rho} \inc t_{max}$, with $\alpha_{t,i,p} \in (0,1]$, to preserve positivity of pressure. According to \eqref{eq:pressure-Jensen} and \eqref{eq:ppLimitDt}, we have
%\begin{equation}
%	p\left(\UM^n_i +  \alpha_{t,i,p} \alpha_{t,i,\rho} \inc t_{max} \R^n_i\right) \geq \left(1-\alpha_{t,i,p}\right)p\left(\UM^n_i\right) + \alpha_{t,i,p} p\left(\UM^n_i + \alpha_{t,i,\rho} \inc t_{max} \R^n_i\right) \geq
%	(1- \eta_t)
%	p\left(\uu_i^n\right),
%\end{equation}
%and thus obtain
and 
\begin{equation}
	\alpha_{\inc,i,p} = \begin{dcases}
		1, & \ \mathrm{if} \ p \left(\UM^{n,s,m}_i + \alpha_{\inc,i,\rho} \inc \uu_i^{n,s,m}\right) \geq  p \left(\UM^{n,s,m}_i\right), \\
		\min\left(1, \dfrac{-\eta_\inc p \left(\UM^{n,s,m}_i\right)}{p \left(\UM^{n,s,m}_i +  \alpha_{\inc,i,\rho} \inc \uu_i^{n,s,m}\right)-  p \left(\UM^{n,s,m}_i\right)} \right), &
		%\mathrm{if} \ p \left(\UM^{n,s,m}_i + \alpha_{\inc,i,\rho} \inc \uu_i^{n,s,m}\right)<  p\left(\UM^{n,s,m}_i\right).
		\mathrm{else}.
	\end{dcases}
\end{equation}
A solution $\uu^{n,s,m+1}_i \in G$ is then computed by
\begin{equation}
	\uu_i^{n,s,m+1} = \uu_i^{n,s,m} + \alpha_{\inc,i} \inc \uu_i^{n,s,m}.
\end{equation}


The positivity-preserving algorithm developed in this paper is based on the scaling limiter \cite{zhang2010positivity}, physical time step limiting, pseudo time step limiting and pseudo time increment correction. The scaling limiter is accuracy-preserving \cite{zhang2010positivity}. It is known that reducing physical time step sizes does not affect the order of accuracy of the time integration. Furthermore, reducing pseudo time step sizes and increments at intermediate steps do not affect the accuracy of the converged solutions of the inner iterations. Therefore, the proposed positivity-preserving approach is accuracy-preserving. Furthermore, it is noted that the proposed positivity-preserving algorithm is applicable to general finite volume methods with implicit dual-time stepping, not limited to finite volume methods with high-order variational reconstruction or high-order ESDIRK time discretization.

%\subsection{Positivity Preserving Limiting of Reconstruction}
%
%Using high order variational reconstruction and WBAP-L2
%accuracy preserving limiter, the reconstructed values
%on quadrature points could still violate positivity-preserving
%even if the mean values $\uu_i$ are positivity-preserving.
%Assuming after reconstruction and shock-capturing limiter,
%the reconstructed quadrature value of conservative quantities are 
%\begin{equation}
%    \U_{i,g} = \uu_{i} + \sum_{l}^{\mathrm{NDOF}}{\U_i^l\varphi_{i,l}(\x_g)}
%\end{equation}
%where $\U_i^l$ are reconstruction coefficients 
%handled with shock-capturing limiter, 
%$g$ is the index of quadrature point and $\x_g$ is the coordinate of 
%the quadrature point. 
%Therefore, in order to preserve the positivity on 
%quadrature points, a simple compression is made as
%\begin{equation}
%    \U_{i,g,pp} = \uu_{i} + \beta_i \sum_{l}^{\mathrm{NDOF}}{\U_i^l\varphi_{i,l}(\x_g)}
%\end{equation}
%so that 
%\begin{equation}
%    \begin{aligned}
%        \rho(\U_{i,g,pp}) &\geq 0,\\
%        e(\U_{i,g,pp}) &\geq 0 ,\ \ \forall i,g
%    \end{aligned}
%\end{equation}
%Following the method for time step limiting, using 
%the convex attribute, 
%a sufficient condition is used to calculate the compression ratio
%\begin{equation}
%    \begin{aligned}
%        \beta_{i} & = \min_g(\beta_{i,g})\\
%        \beta_{i,g} & = \min\left(
%            \frac{\rho(\uu_i^{n,s,m})}{-\rho_-(\inc_{R,g} \uu_i^{n,s,m})},
%            \frac{e(\uu_i^{n,s,m})}{-e_-(\inc_{R,g} \uu_i^{n,s,m})},
%            1
%            \right)^{p_\beta}
%    \end{aligned}
%\end{equation}
%where notation $\inc_{R,g} \uu_i^{n,s,m}$ is 
%\begin{equation}
%    \inc_{R,g} \uu_i^{n,s,m} = \sum_{l}^{\mathrm{NDOF}}{\U_i^{l,n,s,m}\varphi_{i,l}(\x_g)}
%\end{equation}
%which is the reconstructed part of conservative value on the quadrature.
%Usually $p_\beta=1$ but for hard to calculate problems you can set $p_\beta>1$
%to improve stability. The power $p_\beta$ does not affect 
%the asymptote order of $\beta$ when $\beta\approx 1$.
%
%The current research is unable to assert if the positivity-preserving
%limiting of reconstruction affects the overall spacial discretization accuracy.
%In practice, $\beta_i < 1$ only resides in a small portion of 
%the computational domain, therefore the computation is still able to benefit
%form the better accuracy of high order finite volume.

%\subsection{Robust Treatment of Numerical Flux}
%
%The numerical flux on each interface 
%quadrature is dependent on the interface 
%states $\U_{i,g},\U_{j,g}$ and their gradient.
%With reconstruction, even if the positivity-preserving
%limiting of reconstruction coefficients are present and 
%quadrature states are valid physical states, the 
%numerical flux calculated could still be singular
%because the reconstructed values $\U_{i,g},\U_{j,g}$
%may have very small $\rho$ or $e$. 
%Taking LLF flux as an example, a common practice is to use LLF in the form:
%\begin{equation}
%    \label{eq:LLFCommon}
%    \hat{\F} = \F(\U_{i,g}) + \F(\U_{j,g}) - \max(\lambda(\U_{i,g}),\lambda(\U_{j,g}))(\U_{j,g}- \U_{i,g})
%\end{equation}
%where
%\begin{equation}
%    \lambda(\U) = |u_n(\U)| + a(\U)
%\end{equation}
%is the spectral radius at one point, with $u_n(\U)$ being the normal convective speed 
%and $a(\U)$ being the speed of sound.
%When one of $\U_{i,g},\U_{j,g}$ keeps positivity but becomes singular,
%$\lambda(\U)$ could possibly be unbounded.
%A more robust way of calculating LLF flux is using 
%\begin{equation}
%    \hat{\F} = \F(\U_{i,g}) + \F(\U_{j,g}) - \max(\lambda(\uu_i),\lambda(\uu_j))(\U_{j,g}- \U_{i,g})
%\end{equation}
%as the mean values can give more conservative estimations of 
%the eigenvalues.
%
%For more complicated numerical flux terms including those 
%using eigenspace decomposition, viscous terms 
%and temperature dependent properties, 
%the reference value could all be in turn calculated with the mean values $\uu_i, \uu_j$.
%This alternation does not affect the accuracy of the numerical flux 
%as long as the exact flux terms and the difference terms in \eeqref{eq:LLFCommon} use 
%reconstructed values.


\section{Numerical Results}
\label{sec:Results}

The proposed positivity-preserving algorithm is verified using a series of benchmark test cases. The numerical scheme used is the fourth-order accurate implicit finite volume method based on a cubic variational reconstruction and an ESDIRK4 time integration. In the simulations, the relaxation parameters are set as $\eta_t= 0.8$, $\eta_\tau=0.5$ and $\eta_\inc=0.9$, unless otherwise specified. The convergence criterion for inner iteration is that the $L_1$ norm of the pseudo time derivative decreases by three orders of magnitude.

\subsection{Double Rarefaction}

The one-dimensional double rarefaction problem \cite{hu2004kineticDoubleRare} 
is a Riemann problem with the following initial conditions
\begin{equation}
    (\rho,u,p) = \begin{dcases}
    	 (1,-2,0.1),\ \ & x < 0.5,\\
    	(1,2,0.1),\ \ &  \mathrm{else}.
    \end{dcases}
%    \left\{
%        \begin{array}{ll}
%            (1,-2,0.1),\ \ & x < 0.5,\\
%            (1,2,0.1),\ \ &  \mathrm{else}.\\
%        \end{array}
\end{equation}
The flow is inviscid and the ratio of specific heat is $\gamma=1.4$.
In this problem, a vacuum lies in in the middle, adjacent to the ends of two rarefaction waves.
The computational domain is $[0,1]$.
We perform a numerical simulation on a uniform mesh with grid size $\inc x = 1/400$, up to $t=0.1$ using a user-specified time step size $\inc t_{max}= 5 \times 10 ^{-3}$.
In each inner iteration, the CFL number for local pseudo time step $\CFLtau$ is initiated as $0.1$ and increased gradually up to its maximum $10$ at the tenth iteration. 
%The convergence criterion for inner iteration is that the norm of the pseudo time derivative decreases by three orders of magnitude. 
The numerical results are shown in Figure \ref{fig:doubleRare}. It is observed from Figure \ref{fig:doubleRare} that, the numerical solution is essentially oscillation-free and agrees well with the reference solution.

\begin{figure}[htbp]
    \centering
    \begin{subfigure}{0.33\textwidth}
        \includegraphics[width=\textwidth]{pics/PPRobust_DR_R.pdf}
        \caption[]{Density}
    \end{subfigure}\hfill
    \begin{subfigure}{0.33\textwidth}
        \includegraphics[width=\textwidth]{pics/PPRobust_DR_U.pdf}
        \caption[]{Velocity}
    \end{subfigure}\hfill
    \begin{subfigure}{0.33\textwidth}
        \includegraphics[width=\textwidth]{pics/PPRobust_DR_P.pdf}
        \caption[]{Pressure}
    \end{subfigure}
    \caption{Results of the double rarefaction problem at $t=0.1$.}
    \label{fig:doubleRare}
\end{figure}

%\subsection{Le Blanc Shock-Tube}
%
%The Le Blanc shock-tube problem \cite{hu2013positivity} is also a one-dimensional Riemann problem that can be described by the Euler equations. 
%The initial conditions are
%\begin{equation}
%    (\rho,u,p) = \left\{
%        \begin{array}{ll}
%            (1,0,2/3\times10^{-1}),\ \ & x < 3\\
%            (10^{-3},0,2/3\times10^{-10}),\ \ & x \geq 3\\
%        \end{array}
%    \right.
%\end{equation}
%The ratio of specific heat is $\gamma = 5/3$.
%The computational domain is in $[0,9]$, partitioned into $400$ uniform cells.
%In each inner iteration, the CFL number for local pseudo time step $\CFLtau$ is initiated as $0.1$ and increased gradually up to its maximum $10$ at the tenth iteration. 
%%The convergence criterion for inner iteration is that the norm of the pseudo time derivative decreases by three orders of magnitude. 
%Two simulations are performed up to $t=6$, using $600$ time steps with $\inc t_{max} = 0.01$ and  $64$ time steps with $\inc t_{max} = 0.1$, respectively.
%
%\begin{figure}[htbp]
%    \centering
%    \begin{subfigure}{0.33\textwidth}
%        \includegraphics[width=\textwidth]{pics/PPRobust_LB_R.pdf}
%        \caption[]{Density}
%    \end{subfigure}\hfill
%    \begin{subfigure}{0.33\textwidth}
%        \includegraphics[width=\textwidth]{pics/PPRobust_LB_U.pdf}
%        \caption[]{Velocity}
%    \end{subfigure}\hfill
%    \begin{subfigure}{0.33\textwidth}
%        \includegraphics[width=\textwidth]{pics/PPRobust_LB_P.pdf}
%        \caption[]{Pressure}
%    \end{subfigure}
%    \caption{Results of the Le Blanc shock-tube problem at $t=6$}
%    \label{fig:leBlanc}
%\end{figure}
%
%The numerical results of the two simulations are shown in Figure \ref{fig:leBlanc}. It is observed that the numerical solutions with both time step sizes basically agree with the exact solution, and no negative values of pressure and density are detected. 


\subsection{Sedov Blast Wave}

The Sedov blast problem \cite{zhang2012positivity,vilar2016positivity} is a popular case to test positivity-preserving properties of 
numerical schemes. 
The governing equations are the Euler equations and the ratio of specific heat is $\gamma = 1.4$. The computational domain is $[0,1.1]\times[0,1.1]$, partitioned into rectangular cells with grid size $\inc x = \inc y = 1.1/160$. The initial conditions are 
\begin{equation}
    (\rho,u,v,p) = \begin{dcases}
    	(1,0,0,4\times 10 ^{-9}),\ \ & \text{if } x > \inc x \ \mathrm{or} \ y > \inc y, \\
    	(1,0,0,\dfrac{\left(\gamma-1\right)\varepsilon^0}{\inc x\inc y}),\ \ & \text{else},\\
    \end{dcases}
\end{equation}
where $\varepsilon^0$ is the total amount of release energy. By choosing $\varepsilon^0= 2.44816\times 10^5$, the solution consists of a diverging infinite strength shock wave whose front is located at radius $r=1$ at $t=10^{-3}$, with a peak density reaching $6$.
%The final computational time is $t = 1\times10^{-3}$,
%when the exact solution has the shock wave at radius of $1$. 
All boundaries are slip walls.
%Fourth order variational reconstruction finite volume with ESDIRK4
% discretization is used. 
%The relaxation parameters and internal CFL number are the same as 1-D problems.
%Convergence threshold for pseudo time iteration is $10^{-3}$. 
%The convergence criterion for inner iteration is that the norm of the pseudo time derivative decreases by three orders of magnitude. 
%Maximum time step sizes of $\inc t_{max}=10^{-6}$ and $\inc t_{max}=10^{-5}$
%are tested. 
%The computation takes 1026 steps with $\inc t_{max}=10^{-6}$, 
%and 297 steps with $\inc t_{max}=10^{-5}$.
Two simulations are performed up to $t=10^{-3}$, using $1026$ time steps with $\inc t_{max} = 10^{-6}$ and  $297$ time steps with $\inc t_{max} = 10^{-5}$, respectively.


\begin{figure}[htbp]
    \centering
    \begin{subfigure}{0.5\textwidth}
        \includegraphics[width=\textwidth]{pics/PPRobust_SedovDT1.png}
        \caption[]{$\inc t_{max}=10^{-6}$}
    \end{subfigure}\hfill
    \begin{subfigure}{0.5\textwidth}
        \includegraphics[width=\textwidth]{pics/PPRobust_SedovDT10.png}
        \caption[]{$\inc t_{max}=10^{-5}$}
    \end{subfigure}
    \caption{Density contours of the Sedov blast wave problem.}
    \label{fig:sedov}
\end{figure}

The computed density contours at $t=10^{-3}$ are shown in Figure \ref{fig:sedov}. It is observed that the numerical solutions are essentially oscillation-free. To check the solutions in a more intuitive way, we plot both the analytical and numerical density distributions along the $y=x$ line in Figure \ref{fig:sedovLine}.
It is shown in Figure \ref{fig:sedovLine} that, both numerical solutions agree well with the analytic solution,
and the one with the smaller time step size has higher resolution.
%With larger time step size limit $\inc t_{max}=10^{-5}$, the 
%solution appears slightly distorted, but the overall 
%position and strength of the shock front is correct.

\begin{figure}[htbp]
    \centering
    \includegraphics[width=0.6\textwidth]{pics/PPRobust_Sedov.pdf}
    \caption{Density distributions of the Sedov blast wave problem along the diagonal line.}
    \label{fig:sedovLine}
\end{figure}






\subsection{Mach 2000 Jet}

The Mach 2000 jet problem \cite{zhang2010positivity} is a challenging case as it has extremely strong discontinuities.
Following the practice of \cite{huang2024general}, 
the current paper studies two cases with $Re=\infty$ and $Re=100$ to demonstrate the capability of the proposed positivity-preserving algorithm to deal with inviscid and viscous compressible flows, respectively.
The governing equations are the Navier-Stokes equations with $\gamma=5/3$.
The computational domain is $[0,1]\times[-0.25,0.25]$.
Far-field boundary conditions are imposed on the upper, lower and right boundaries. The following inflow conditions are imposed on the left boundary:
\begin{equation}
    (\rho,u,v,p) = \left\{
        \begin{array}{ll}
            (5,800,0,0,4127),\ \ & \text{if } -0.05 \leq y \leq 0.05, \\
            (0.5,0,0,0.4127),\ \ & \text{else}.\\
        \end{array}
    \right.
\end{equation}
The initial conditions are $(\rho,u,v,p)=(0.5,0,0,0.4127)$.
The numerical simulations are performed on a uniform rectangular mesh with $800 \times 800$ cells, up to $t=10^{-3}$ using a user-specified time step size $\inc t_{max} = 1\times10^{-6}$. 
The relaxation parameters are  $\eta_t= 0.6$, $\eta_\tau=0.5$, $\eta_\inc=0.9$.
%$\CFLtau$ is initially $0.1$ at the beginning of each stage,
%and reaches $2.0$ at the 10th pseudo time iteration. 
In each inner iteration, the CFL number for local pseudo time step $\CFLtau$ is initiated as $0.1$ and increased gradually up to its maximum $2$ at the tenth iteration. 

The computed density and pressure contours at $t=10^{-3}$ in logarithmic scales are shown in Figure \ref{fig:M2000_ReInf_R}, \ref{fig:M2000_Re1E2_R}, \ref{fig:M2000_ReInf_P} and \ref{fig:M2000_Re1E2_P}.
It is observed from these figures that the numerical results are essentially non-oscillatory, and small-scale structures are smeared out in the $Re=100$ case due to the viscous effect.

\begin{figure}[htbp]
    \centering
    \includegraphics[trim={5px 0 5px 0},clip,width=0.9\textwidth]{pics/PPRobust_M2000_ReInf_R.png}
    \caption{$\log_{10}(\rho)$ of $Re=\infty$ Mach 2000 jet, 40 contour lines from -1.5 to 1.5.}
    \label{fig:M2000_ReInf_R}
\end{figure}

\begin{figure}[htbp]
	\centering
	\includegraphics[trim={5px 0 5px 0},clip,width=0.9\textwidth]{pics/PPRobust_M2000_Re1E2_R.png}
	\caption{$\log_{10}(\rho)$ of $Re=100$ Mach 2000 jet, 40 contour lines from -1.5 to 1.5.}
	\label{fig:M2000_Re1E2_R}
\end{figure}

\begin{figure}[htbp]
	\centering
	\includegraphics[trim={5px 0 5px 0},clip,width=0.9\textwidth]{pics/PPRobust_M2000_ReInf_P.png}
	\caption{$\log_{10}(p)$ of $Re=\infty$ Mach 2000 jet, 40 contour lines from -1.4 to 5.3.}
	\label{fig:M2000_ReInf_P}
\end{figure}

\begin{figure}[htbp]
    \centering
    \includegraphics[trim={5px 0 5px 0},clip,width=0.9\textwidth]{pics/PPRobust_M2000_Re1E2_P.png}
    \caption{$\log_{10}(p)$ of $Re=100$ Mach 2000 jet, 40 contour lines from -1.4 to 5.3.}
    \label{fig:M2000_Re1E2_P}
\end{figure}

\subsection{Shock Diffraction}

The shock diffraction problem \cite{zhang2010positivity}, in which a shock passes a backward facing corner, is used to test the positivity-preserving capability of the proposed algorithm. 
The governing equations are the Euler equations with $\gamma = 1.4$.
The computational domain is the union of $[0,1]\times[6,11]$ and $[1,13]\times[0,11]$.
The initial condition is a Mach number $5.09$ shock located at $x=0.5$, moving into undisturbed air ahead of the shock at a state $(\rho, u, v, p)=(1.4,0,0,1)$.
%{\color{red}The boundary conditions are inflow at $x = 0, \ 6\leq y \leq 11$, outflow at $x = 13, \ 0\leq y \leq 11$, $1 \leq x \leq 13, \ y = 0$ and $0 \leq x \leq 13, \ y = 11$, and non-slip at the walls $0\leq x \leq 1, \ y = 6$ and $x = 1, \ 0\leq y \leq 6$.}
Inflow and slip solid wall boundary conditions are imposed on the left side $x = 0$, $6\leq y \leq 11$ and other boundaries, respectively.
%The boundary at $x=0$ is set to match the left state of the shock, and 
%all other boundaries are set as inviscid wall.
The simulation is performed on a rectangular mesh with grid size $\inc x = \inc y = 1/80$, up to $t=2.3$ using a user-specified time step size $\inc t_{max}=1\times 10^{-3}$.
The CFL number for pseudo-time step is $\CFLtau = 1$. 

\begin{figure}[htbp]
    \centering
    \begin{subfigure}{0.5\textwidth}
        \includegraphics[width=\textwidth]{pics/PPRobust_Corner_R.png}
        \caption[]{$\log_{10}(\rho)$, 40 contour lines from -0.8 to 0.8}
    \end{subfigure}\hfill
    \begin{subfigure}{0.5\textwidth}
        \includegraphics[width=\textwidth]{pics/PPRobust_Corner_P.png}
        \caption[]{$\log_{10}(p)$, 40 contour lines from -0.6 to 1.6}
    \end{subfigure}
    \caption{Density and pressure contours of the shock diffraction problem.}
    \label{fig:Corner}
\end{figure}

The computed density and pressure contours in logarithmic scales at $t=2.3$ are shown in Figure \ref{fig:Corner}. It is observed from Figure \ref{fig:Corner} that, the numerical solutions are essentially oscillation-free and small-scale shock structures are well resolved.


\subsection{Shock Reflection and Diffraction Around a Wedge}

The shock reflection and diffraction problem \cite{zhang2017positivity} is similar to the shock diffraction problem, with the corner replaced by a wedge. The computational domain is $[0,3]\times[0,2]$ with a $30^\circ$ wedge placed at $x\in[0.2,1.2]$, while the 
tip of the wedge is at $(x,y) = (1.2, 1/\sqrt{3})$. 
Initially, there is a Mach $10$ shock at $x=0.2$ moving into undisturbed air at a state $(\rho, u, v, p) = (1.4,0,0,1)$. The $x=0$ boundary on 
the left matches the downstream state of the shock, and all other boundaries 
are inviscid walls.
Before the diffraction of the shock, the development
of the shock structure is identical with that in the Mach $10$ double Mach reflection problem \cite{woodward1984dmr}. Similar to the Mach $2000$ jet example, $\Re=\infty$ and $\Re=100$ cases 
are investigated. A triangular mesh with grid size $h = 1/160$, shown in Figure \ref{fig:wedgeMesh}, is used in the simulations.
The simulations are performed up to $t=0.245$ using
a user-specified time step size $\inc t_{max}= 2\times10^{-4}$.
In each inner iteration, the CFL number for local pseudo time step $\CFLtau$ is initiated as $0.01$ and increased gradually up to its maximum $1$ at the tenth iteration. 

\begin{figure}[htbp]
    \centering
    \includegraphics[trim={5px 0 5px 0},clip,width=0.6\textwidth]{pics/PPRobust_WedgeMesh.png}
    \caption{Mesh used in the shock reflection and diffraction problem.}
    \label{fig:wedgeMesh}
\end{figure}

The computed density and pressure contours are shown in Figure \ref{fig:wedgeReInf} and \ref{fig:wedgeRe1E2}. 
It is observed that the numerical solutions are essentially non-oscillatory. The Kelvin-Helmholtz instability induced in the shear layer is well resolved in Figure \ref{fig:wedgeReInf} for the $\Re=\infty$ case. While in Figure  \ref{fig:wedgeRe1E2} for the $\Re=100$ case, the small-scale flow structures around the shear layer and the wedge are smeared  due to the existence of physical viscosity.

\begin{figure}[htbp]
    \centering
    \begin{subfigure}{0.5\textwidth}
        \includegraphics[width=\textwidth]{pics/PPRobust_WedgeReInf_R.png}
        \caption[]{$\rho$, 60 contour lines from 0 to 22}
    \end{subfigure}\hfill
    \begin{subfigure}{0.5\textwidth}
        \includegraphics[width=\textwidth]{pics/PPRobust_WedgeReInf_P.png}
        \caption[]{$\log_{10}(p)$, 60 contour lines from -1.7 to 2.7}
    \end{subfigure}
    \caption{Density and pressure contours for the shock reflection and diffraction problem, $\Re=\infty$.}
    \label{fig:wedgeReInf}
\end{figure}

\begin{figure}[htbp]
    \centering
    \begin{subfigure}{0.5\textwidth}
        \includegraphics[width=\textwidth]{pics/PPRobust_WedgeRe1E2_R.png}
        \caption[]{$\rho$, 60 contour lines from 0 to 22}
    \end{subfigure}\hfill
    \begin{subfigure}{0.5\textwidth}
        \includegraphics[width=\textwidth]{pics/PPRobust_WedgeRe1E2_P.png}
        \caption[]{$\log_{10}(p)$, 60 contour lines from -1.7 to 2.7}
    \end{subfigure}
    \caption{Density and pressure contours for the shock reflection and diffraction problem, $\Re=100$.}
    \label{fig:wedgeRe1E2}
\end{figure}


\section{Conclusions}
\label{sec:Conclusions}

This paper presents a positivity-preserving algorithm for finite volume methods using dual-time stepping for compressible flow simulations. In the positivity-preserving algorithm, admissible cell-averages are obtained by limiting solution changes, which is accomplished by controlling physical and pseudo time step sizes. To overcome the difficulty of unknown solution changes, we propose to use explicit time discretizations to obtain efficient estimations of future states. The allowable time step sizes are determined by limiting the relative solution changes. Given positive cell-averages, admissible reconstruction polynomials can be obtained by applying a positivity-preserving scaling limiter. The proposed positivity-preserving algorithm is accuracy-preserving. Numerical results for a series of benchmark test cases demonstrate the robustness and high resolution of the positivity-preserving implicit high-order finite volume method. 



%%%% BIBLIOGRAPHY
\bibspacing=\dimen 100 
\bibliographystyle{unsrt} 
\bibliography{refs}

\end{document}
